\documentclass[a4paper,12pt]{article}
\usepackage[utf8]{inputenc}
\usepackage[ngerman]{babel}
\usepackage{amsmath,amssymb}
\usepackage{geometry}
\geometry{margin=2.5cm}

\title{Aufgabe 7(b) -- Gauß- und Gauß-Jordan-Elimination}
\author{}
\date{}

\begin{document}
\maketitle

\section*{Gleichungssystem}

\[
\begin{cases}
4x + 7y = 1 \\
y - 2z = -5 \\
2x + y = -3 \\
x + y + z = 2
\end{cases}
\]

\textbf{Hinweis:} Dieses System unterscheidet sich von 7(a) nur in der rechten Seite der ersten Gleichung ($1$ statt $-1$).

\section*{Erweiterte Koeffizientenmatrix}

\[
\left(\begin{array}{ccc|c}
4 & 7 & 0 & 1 \\
0 & 1 & -2 & -5 \\
2 & 1 & 0 & -3 \\
1 & 1 & 1 & 2
\end{array}\right)
\]

\section*{Gauß-Elimination (Stufenform)}

\textbf{Schritt 1:} $Z_1 \leftrightarrow Z_4$

\[
\left(\begin{array}{ccc|c}
1 & 1 & 1 & 2 \\
0 & 1 & -2 & -5 \\
2 & 1 & 0 & -3 \\
4 & 7 & 0 & 1
\end{array}\right)
\]

\textbf{Schritt 2:} $Z_3 \leftarrow Z_3 - 2Z_1$ und $Z_4 \leftarrow Z_4 - 4Z_1$

\[
\left(\begin{array}{ccc|c}
1 & 1 & 1 & 2 \\
0 & 1 & -2 & -5 \\
0 & -1 & -2 & -7 \\
0 & 3 & -4 & -7
\end{array}\right)
\]

\textbf{Schritt 3:} $Z_3 \leftarrow Z_3 + Z_2$ und $Z_4 \leftarrow Z_4 - 3Z_2$

\[
\left(\begin{array}{ccc|c}
1 & 1 & 1 & 2 \\
0 & 1 & -2 & -5 \\
0 & 0 & -4 & -12 \\
0 & 0 & 2 & 8
\end{array}\right)
\]

\textbf{Schritt 4:} $Z_4 \leftarrow 2Z_4 + Z_3$

\[
\left(\begin{array}{ccc|c}
1 & 1 & 1 & 2 \\
0 & 1 & -2 & -5 \\
0 & 0 & -4 & -12 \\
0 & 0 & 0 & 4
\end{array}\right)
\]

\section*{Ergebnis}

Die letzte Zeile lautet:
\[
0x + 0y + 0z = 4
\]
Das ist ein \textbf{Widerspruch} ($0 = 4$).

\[
\boxed{\text{Das Gleichungssystem hat \textbf{keine Lösung}.}}
\]

\subsection*{Erklärung}

Das System ist \textbf{inkonsistent}. Die vier Gleichungen widersprechen sich gegenseitig.
Im Vergleich zu 7(a) wurde nur die rechte Seite der ersten Gleichung von $-1$ auf $1$ geändert --
das reicht aus, um das System unlösbar zu machen.

Da das System inkonsistent ist, wird \textbf{keine} Gauß-Jordan-Elimination durchgeführt
(diese ist nur für konsistente Systeme gefordert).

\end{document}
